\documentclass[11pt]{article}
\usepackage[a4paper,margin=1in]{geometry}
\usepackage{amsmath,amssymb,mathtools,bm}
\usepackage{enumitem,booktabs,array}
\usepackage{tcolorbox}
\usepackage{hyperref}
\usepackage[protrusion=true,expansion=false]{microtype}
\hypersetup{colorlinks=true,linkcolor=blue,urlcolor=blue,citecolor=blue}
\setlist[itemize]{topsep=2pt,itemsep=2pt}
\setlist[enumerate]{topsep=2pt,itemsep=2pt}
\newtcolorbox{keybox}{colback=blue!3!white,colframe=blue!40!black,boxrule=0.4pt,arc=1mm}
\newtcolorbox{alertbox}{colback=red!3!white,colframe=red!40!black,boxrule=0.4pt,arc=1mm}
\newtcolorbox{answerbox}{colback=green!3!white,colframe=green!40!black,boxrule=0.4pt,arc=1mm}
\newtcolorbox{drillbox}{colback=yellow!5!white,colframe=orange!60!black,boxrule=0.4pt,arc=1mm}
\newcommand{\EV}{\mathbb{E}}
\newcommand{\Var}{\mathrm{Var}}
\newcommand{\Cov}{\mathrm{Cov}}
\newcommand{\blank}[1]{\underline{\hspace{#1}}}

\title{E12-01: Harford, Wang \& Zhang (2017, RFS) \\
\large ``Foreign Cash: Taxes, Internal Capital Markets, and Agency Problems'' --- Active Recall Workbook}
\author{Empirical Corporate Finance II --- Exam Prep}
\date{\today}
\begin{document}\maketitle

\begin{keybox}
\textbf{Paper in one sentence.} HWZ show that US multinationals accumulate cash in foreign subsidiaries because of the repatriation tax, that this cash is trapped from the US parent's perspective, and that high foreign cash is associated with lower firm value and more value-destroying foreign acquisitions --- evidence that the tax-driven foreign-cash trap creates internal-capital-market agency problems.

\textbf{Placement.} Key paper on US multinationals' cash holdings and the repatriation-tax distortion. Informs the 2017 TCJA debate on deemed-repatriation.
\end{keybox}

\section*{Round 1: Complete Walk-Through}

\subsection*{1.1 Institutional setting}
Pre-TCJA (pre-2017), US taxed worldwide income with a credit for foreign taxes paid; firms deferred US tax on foreign earnings until repatriation. This created strong incentives to accumulate cash offshore.

\subsection*{1.2 Data}
\begin{itemize}
\item S\&P 1500 multinational firms 1997--2013.
\item 10-K disclosures of permanently-reinvested foreign earnings (PRE); cash held abroad.
\item Foreign acquisitions from SDC.
\end{itemize}

\subsection*{1.3 Empirical design}
\begin{itemize}
\item Cross-sectional regressions of firm value (Tobin's $q$) on foreign-cash ratios.
\item Event studies of foreign-M\&A announcement returns conditional on foreign-cash holdings.
\item DiD around changes in effective repatriation-tax rates (e.g., 2004 HIA holiday).
\end{itemize}

\subsection*{1.4 Headline results}
\begin{itemize}
\item Foreign cash grew from $\sim$5\% to 15--20\% of assets at affected multinationals.
\item High-foreign-cash firms earn lower CARs on foreign acquisitions ($-$1 to $-$2 pp).
\item These deals are disproportionately value-destroying and driven by availability of trapped cash.
\item Firms substantially increase US debt issuance to fund US operations despite having cash abroad.
\end{itemize}

\subsection*{1.5 Mechanism}
\begin{itemize}
\item Repatriation tax creates a \emph{trapped-cash} internal-capital-market friction.
\item Managers deploy trapped cash in foreign investments to avoid the tax, choosing marginal/low-NPV deals.
\item Effective cost of capital for foreign investment is artificially low.
\end{itemize}

\subsection*{1.6 Policy implications}
\begin{itemize}
\item Supports the TCJA's shift to a near-territorial regime.
\item Deemed-repatriation provisions in TCJA directly address the trapped-cash problem.
\item Provides ex-ante evidence of the inefficiency the reform was designed to correct.
\end{itemize}

\subsection*{1.7 Caveats}
\begin{alertbox}
\begin{itemize}
\item PRE reporting is subject to disclosure rules; measurement error possible.
\item Foreign-M\&A selection: firms with trapped cash may naturally face different opportunities.
\item Dynamic tax-planning confounds cross-sectional inference.
\end{itemize}
\end{alertbox}

\section*{Round 2: Level 1 Fill-in}
\begin{enumerate}
\item Institutional friction: US \blank{2cm} tax on foreign earnings.
\item Foreign cash growth: $5\%\to$\blank{1cm}--\blank{1cm}\% of assets.
\item Effect on foreign-M\&A CARs: $-$\blank{1cm} to $-$\blank{1cm} pp.
\item Mechanism: \blank{2cm}-cash internal-capital-market friction.
\item Policy response: \blank{2cm} 2017 (TCJA).
\end{enumerate}

\section*{Round 3: Level 2 Prompts}
\begin{enumerate}
\item Describe the pre-TCJA repatriation-tax setting.
\item Report main findings on foreign cash and foreign-M\&A value destruction.
\item Explain the trapped-cash mechanism.
\item Discuss 2004 HIA and 2017 TCJA reforms.
\item Evaluate endogeneity concerns.
\end{enumerate}

\section*{Round 4: Minimal Scaffolding}
From a blank page: (i) setting, (ii) data, (iii) findings, (iv) mechanism, (v) policy.

\section*{Round 5: Blank Page}
Essay: trapped foreign cash and US multinational capital allocation.

\vspace{6pt}
\begin{drillbox}
\textbf{Speed Drill.} (1) Tax friction. (2) Foreign-cash growth. (3) M\&A CAR effect. (4) Mechanism. (5) Policy reform.
\end{drillbox}

\section*{Quick Reference \& Answer Key}
\begin{answerbox}
\begin{itemize}
\item \textbf{Friction.} US repatriation tax traps foreign cash.
\item \textbf{Effect.} Higher foreign cash $\to$ worse foreign M\&A.
\item \textbf{Mechanism.} Agency + internal-capital distortion.
\item \textbf{Reform.} TCJA 2017 territorial regime.
\end{itemize}
\end{answerbox}

\begin{keybox}
\textbf{30-second exam pitch.} Harford, Wang \& Zhang (2017) study how the US pre-TCJA worldwide corporate-tax regime --- specifically the repatriation tax on foreign earnings --- shaped multinationals' cash management and capital allocation. Using S\&P 1500 multinational 10-K disclosures of permanently-reinvested foreign earnings 1997--2013 and SDC foreign-M\&A data, they show foreign cash grew from $\sim$5\% to 15--20\% of assets, high-foreign-cash firms earn $-$1 to $-$2 pp lower announcement CARs on foreign acquisitions, and US operations are funded with new debt despite offshore cash. Trapped cash creates an internal-capital-market agency friction: managers deploy it in marginal foreign deals to avoid the repatriation tax. The paper provides the clearest ex-ante evidence of the inefficiency that the 2017 TCJA's deemed-repatriation and territorial-tax reforms were designed to correct.
\end{keybox}

\end{document}
